\chapter{El simulador QRS/ST}
\label{chap:simulador}

\section{Objetivo del dominio sintético}

Los datos sintéticos permiten estudiar la eficiencia con pocos ejemplos sin depender de una colección clínica extensa. En este trabajo, el simulador funciona como un entorno controlado. Cada latido se construye a partir de componentes interpretables, se somete a reglas explícitas y se renderiza como una imagen. Tanto la CNN como el modelo de visión y lenguaje reciben esa representación visual.

El objetivo no consiste en reproducir la fisiología cardíaca completa. Para ello habría que modelar las doce derivaciones, los cambios de frecuencia y la interacción entre ciclos. También serían necesarias la conducción auricular y ventricular, el ruido de adquisición, la variación anatómica y el efecto del dispositivo. El sistema implementado se limita a un latido aislado con apariencia de derivación precordial derecha. Esta simplificación hace posible parametrizar los tres hallazgos visuales de interés para el experimento.

El uso de datos sintéticos cumple cuatro funciones:

\begin{enumerate}
\item Generar etiquetas sin depender de anotación manual para cada imagen.
\item Separar la forma de una onda de factores externos como fabricante, hospital o estilo de documento.
\item Construir conjuntos de ejemplos con tamaños exactos y semillas conocidas.
\item Reservar una familia generadora compuesta para estudiar transferencia.
\end{enumerate}

Esta capacidad de control tiene como contrapartida una brecha entre los dominios sintético y real. El simulador permite comprobar si los modelos aprenden una tarea definida mediante reglas, pero no anticipa el rendimiento en ECG clínicos. Los resultados se interpretan teniendo presente esta diferencia.

\section{Representación matemática}

Cada latido se expresa como la suma de diez átomos gaussianos. En el átomo \(i\), la amplitud \(a_i\), el centro temporal \(\mu_i\) y la anchura \(\sigma_i\) definen la función

\begin{equation}
g_i(t) =
a_i
\exp
\left(
-\frac{(t-\mu_i)^2}{2\sigma_i^2}
\right).
\label{eq:gaussian_atom_sim}
\end{equation}

La señal completa es

\begin{equation}
x(t) =
\sum_{i=1}^{10} g_i(t).
\label{eq:sum_gaussians}
\end{equation}

El signo de \(a_i\) determina la dirección de la deflexión y su valor absoluto, la magnitud. El parámetro \(\mu_i\) sitúa el componente en el eje temporal, mientras que \(\sigma_i\) controla su extensión. Al sumar los átomos se obtienen transiciones suaves y sin discontinuidades.

El uso de gaussianas parte de trabajos previos que representan las ondas electrocardiográficas mediante componentes paramétricos \cite{sayadi2010}. Este simulador es más sencillo que un modelo electrofisiológico dinámico, pero permite relacionar cada modificación con una región concreta de la señal.

\begin{table}[H]
\centering
\small
\begin{tabularx}{\textwidth}{p{2.2cm} p{3.8cm} X}
\toprule
\textbf{Átomo} & \textbf{Región principal} & \textbf{Función dentro del generador} \\
\midrule
\texttt{P} & Onda P & Añade una deflexión auricular positiva anterior al QRS \\
\texttt{Q} & Inicio del QRS & Permite una pequeña deflexión negativa, aunque la configuración actual la mantiene anulada \\
\texttt{R} & QRS inicial & Produce el primer máximo positivo ventricular \\
\texttt{S} & QRS central & Forma la deflexión negativa posterior a R \\
\texttt{R\_prime} & QRS terminal & Introduce un segundo máximo positivo y permite una forma RSR prima \\
\texttt{J\_elev} & Punto J & Controla la transición entre el final del QRS y el segmento ST \\
\texttt{ST1} & ST inicial & Modela elevación temprana y la parte descendente de una forma cóncava \\
\texttt{ST2} & ST tardío & Controla la continuidad hacia la onda T \\
\texttt{T1} & Onda T principal & Define la polaridad y la magnitud principal de la repolarización \\
\texttt{T2} & Cola de la onda T & Suaviza y prolonga la repolarización \\
\bottomrule
\end{tabularx}
\caption{Átomos utilizados para construir un latido sintético}
\label{tab:atomos}
\end{table}

La Figura~\ref{fig:gaussian_generator} muestra cómo se combinan los átomos de la Tabla~\ref{tab:atomos}.

\begin{figure}[H]
  \centering
  \safeincludegraphics[width=\textwidth]{results/gaussian_generator.drawio.png}
  \caption{Generador gaussiano de un latido sintético. Los parámetros de cada átomo fijan su amplitud, posición y anchura antes de sumar las regiones P, QRS, ST y T.}
  \label{fig:gaussian_generator}
\end{figure}

\section{Muestreo temporal}

La configuración final utiliza una frecuencia de 500 Hz y una duración de 900 ms. El número de muestras se calcula como

\begin{equation}
N =
\operatorname{round}
\left(
\frac{D}{1000}f_s
\right),
\end{equation}

donde \(D\) es la duración en milisegundos y \(f_s\) la frecuencia de muestreo. Con los valores finales se obtienen 450 muestras. El vector temporal es

\begin{equation}
t_n =
\frac{1000n}{f_s},
\qquad
n \in \{0,\ldots,N-1\}.
\end{equation}

Los rangos de repolarización se definieron para una duración de referencia de 900 ms. Cuando cambia la duración del latido, los centros posteriores a 390 ms se adaptan mediante

\begin{equation}
\mu_i^{(D)} =
390 +
\left(
\mu_i^{(900)}-390
\right)
\frac{D-390}{900-390}.
\label{eq:time_scaling}
\end{equation}

Las anchuras de los componentes tardíos se escalan con el mismo factor y se impone un mínimo para evitar gaussianas degeneradas. Los componentes anteriores al punto de pivote permanecen fijos. Así, el QRS conserva su posición, mientras que ST y T se comprimen o expanden según la duración disponible.

\section{Rangos de parámetros}

La generación comienza con una configuración común que asigna a cada átomo intervalos para \(a\), \(\mu\) y \(\sigma\). A continuación se sustituyen los intervalos propios de la familia seleccionada. Aunque el muestreo dentro de cada rango es uniforme, la distribución de las señales aceptadas no lo es: las reglas morfológicas y los validadores descartan parte de las propuestas.

Estos rangos pertenecen al generador visual y no representan límites clínicos. Por ejemplo, el átomo \texttt{R\_prime} mantiene una amplitud casi nula en la familia normal y admite valores mayores en RBBB. De forma análoga, \texttt{J\_elev}, \texttt{ST1} y \texttt{ST2} permanecen cerca de cero en las señales normales y adoptan amplitudes positivas en la familia con elevación del ST.

Los rangos favorecen una morfología, mientras que los validadores deciden si la señal resultante satisface la definición de cada etiqueta. Mantener separadas ambas funciones evita que un cambio en la generación modifique de forma inadvertida el significado de las clases.

\section{Familias generadoras}

El simulador parte de cinco familias. Cada una selecciona rangos de muestreo y reglas de coherencia, pero no asigna directamente el vector objetivo.

\subsection{Familia normal}

La familia \texttt{NORMAL} reduce la amplitud de R prima, mantiene las regiones J y ST próximas a la línea base y genera ondas T positivas. Además, rechaza cualquier señal que active alguno de los tres validadores de hallazgos. Esta comprobación evita que una combinación obtenida por azar dentro de los rangos normales supere uno de sus umbrales.

\subsection{Morfología RBBB}

La familia \texttt{RBBB} favorece una segunda deflexión positiva después de S. El centro de R prima se sitúa al menos 24 ms después del centro de S. Dentro de los límites permitidos, su amplitud debe alcanzar como mínimo el 80 por ciento de la amplitud de R. También se le asigna una anchura mayor que a la primera R.

El propósito es obtener una forma RSR prima visible, no reproducir todos los criterios clínicos del bloqueo completo de rama derecha. Estos dependen, entre otros aspectos, de la duración global del QRS y de su distribución entre derivaciones \cite{surawicz2009}. Por este motivo, el texto interpreta la etiqueta como una morfología compatible.

\subsection{Elevación de ST}

La familia \texttt{ST\_ELEVATION} asigna amplitudes positivas al punto J y a los dos componentes del ST. La primera sección del segmento se ensancha para que no adopte la forma de un pico estrecho confundible con una deflexión del QRS. El validador también exige actividad ventricular positiva y negativa suficiente, lo que permite descartar señales planas desplazadas hacia valores altos.

\subsection{Inversión de T}

La familia \texttt{T\_WAVE\_INVERSION} hace negativos los dos componentes de T y los separa de la región ST. Puede generar, además, una R prima moderada y una pequeña elevación de J, lo que permite coocurrencias con RBBB y elevación del ST. Para aceptar la señal, la deflexión negativa debe aparecer tarde, alcanzar una profundidad mínima y mantenerse durante un intervalo suficiente.

\subsection{Familia Brugada}

La familia \texttt{BRUGADA} representa de forma esquemática una morfología coved o de tipo 1. Combina elevación de J, un ST inicial positivo, un descenso posterior y la inversión de T. La amplitud de R prima es menor que en la familia RBBB, aunque puede conservar una alteración terminal del QRS.

El nombre de la familia no forma parte del vector objetivo. Una señal aceptada por el validador interno de Brugada se evalúa después con los validadores independientes de RBBB, elevación del ST e inversión de T. Puede recibir, por tanto, una combinación parcial de hallazgos. La familia se utiliza así como prueba de composición, sin introducir una etiqueta clínica que el simulador no permite justificar.

\begin{table}[H]
\centering
\small
\begin{tabularx}{\textwidth}{p{3.1cm} p{4.4cm} X}
\toprule
\textbf{Familia} & \textbf{Modificación dominante} & \textbf{Interpretación experimental} \\
\midrule
\texttt{NORMAL} & ST próximo a cero, T positiva y R prima mínima & Referencia sin hallazgos del conjunto \\
\texttt{RBBB} & Segundo pico positivo tardío y ancho & Componente terminal del QRS \\
\texttt{ST} elevado & Punto J y ST positivos & Componente de elevación posterior al QRS \\
Inversión de T & Repolarización tardía negativa & Componente de inversión de T \\
\texttt{BRUGADA} & J elevado, ST descendente y T negativa & Familia compuesta reservable para transferencia \\
\bottomrule
\end{tabularx}
\caption{Relación entre familias generadoras y rasgos favorecidos}
\label{tab:familias}
\end{table}

\section{Reglas de coherencia morfológica}

El muestreo independiente de los átomos puede alterar el orden temporal de las ondas. Para evitarlo, antes de sumar las gaussianas se imponen varias relaciones. Q precede a R; S aparece después de R y tiene mayor anchura; J se sitúa después de S. Los dos componentes del ST se ordenan tras J, y T1 y T2 ocupan la región final.

Estas relaciones no convierten la señal en una representación fisiológica. Su función es mantener una geometría mínima que preserve la interpretación de los parámetros. Las reglas propias de cada familia añaden restricciones como la posición tardía de R prima o el signo negativo de T.

Los cambios en una regla pueden afectar a más de un componente. Desplazar T1, por ejemplo, modifica el resultado del validador de inversión y también puede alterar la media tardía del ST. Por ello, cualquier ajuste requiere volver a generar ejemplos y revisar tanto las frecuencias de las etiquetas como las tasas de aceptación.

\section{Muestreo con rechazo}

La generación utiliza muestreo con rechazo. Para una familia \(c\), el procedimiento consta de los siguientes pasos:

\begin{enumerate}
\item Combinar los rangos base con los rangos específicos de \(c\).
\item Adaptar los componentes tardíos a la duración solicitada.
\item Muestrear los parámetros de los diez átomos.
\item Aplicar las reglas de orden y coherencia.
\item Sumar los átomos sobre el vector temporal.
\item Evaluar el validador correspondiente a \(c\).
\item Aceptar la señal si supera las reglas o volver al paso de muestreo.
\end{enumerate}

El generador admite un máximo de 400 intentos. Si ninguno produce una señal válida, devuelve un error junto con las últimas métricas calculadas. Los metadatos de cada muestra registran el número de intentos, lo que ayuda a detectar rangos poco compatibles con los validadores.

Sea \(q_c\) la probabilidad de aceptación de una propuesta de la familia \(c\). El número esperado de intentos antes de una aceptación es aproximadamente

\begin{equation}
\mathbb{E}[A_c] =
\frac{1}{q_c}.
\end{equation}

Un número elevado de intentos no invalida por sí mismo las imágenes aceptadas, pero indica que gran parte del espacio paramétrico queda fuera de la definición de la etiqueta. Por ello, la tasa de aceptación forma parte de la revisión del generador.

\section{Validadores y semántica de las etiquetas}

Los validadores analizan ventanas temporales de la señal completa. Las situadas después de 390 ms se adaptan mediante la ecuación~\ref{eq:time_scaling}. Las amplitudes se expresan en mV y las duraciones, en ms.

\begin{table}[H]
\centering
\scriptsize
\begin{tabularx}{\textwidth}{p{2.25cm} X}
\toprule
\textbf{Validador} & \textbf{Condiciones necesarias} \\
\midrule
RBBB &
Dos máximos positivos entre 240 y 470 ms, separación mínima de 28 ms, primer máximo superior a 0,12 mV, segundo máximo superior a 0,20 mV, mínimo intermedio inferior a \(-0,05\) mV, segundo máximo posterior a 345 ms y anchura terminal superior a 38 ms por encima de 0,10 mV \\
Elevación de ST &
Máximo del QRS superior a 0,28 mV, mínimo del QRS inferior a \(-0,14\) mV, media de J superior a 0,09 mV y media del ST superior a 0,09 mV \\
Inversión de T &
Máximo del QRS superior a 0,22 mV, máximo de R prima superior a 0,08 mV, media del ST superior a \(-0,03\) mV, mínimo de T inferior a \(-0,12\) mV, mínimo posterior al umbral temporal equivalente a 640 ms y más de 45 ms por debajo de \(-0,05\) mV \\
Normal &
Máximo del QRS superior a 0,22 mV, mínimo inferior a \(-0,10\) mV, media de ST entre \(-0,04\) y 0,05 mV, máximo de T superior a 0,12 mV, R prima inferior a 0,10 mV y rechazo de los tres hallazgos \\
Brugada &
Máximo del QRS superior a 0,18 mV, mínimo inferior a \(-0,16\) mV, media de J superior a 0,09 mV, media de ST superior a 0,08 mV, ST temprano al menos 0,03 mV por encima del ST tardío y mínimo de T inferior a \(-0,08\) mV \\
\bottomrule
\end{tabularx}
\caption{Criterios de los validadores sintéticos}
\label{tab:validadores}
\end{table}

Los umbrales de la Tabla~\ref{tab:validadores} separan morfologías dentro del simulador. No son criterios clínicos ni deben aplicarse a ECG reales, cuya interpretación requiere calibración, derivaciones identificadas, velocidad de papel conocida y contexto clínico. Los estándares revisados en el capítulo~\ref{chap:fundamentos} solo orientan la forma general de las señales sintéticas.

Las tres etiquetas finales se calculan de forma independiente. Una vez aceptada la señal, se ejecutan los validadores de RBBB, elevación del ST y TWI para construir un vector binario. Las coocurrencias resultantes no tienen por qué coincidir con la familia de origen: la etiqueta describe la señal generada, no el mecanismo utilizado para proponerla.

\section{Renderizado de las señales}

Cada latido se representa en un lienzo cuadrado de 896 por 896 píxeles y 130 puntos por pulgada. En el eje horizontal, la cuadrícula menor marca intervalos de 40 ms y la mayor, de 200 ms. En el eje vertical, las divisiones corresponden a 0,1 y 0,5 mV, respectivamente. La señal se dibuja en negro con un grosor de dos puntos y la cuadrícula, en dos tonos de rosa.

El renderizado conserva los títulos de los ejes y los valores numéricos. Aunque aportan información de escala, podrían actuar como atajos si los límites variasen entre muestras. Para evitarlo, todas las imágenes del experimento final comparten los mismos límites horizontales y verticales. También se revisó que ninguna familia presentase un encuadre, un margen o un patrón de marcas propio.

No se aplican reflejos horizontales, pues invertir el eje temporal alteraría el orden de las ondas, ni verticales, porque intercambiarían elevación y depresión. Las transformaciones admisibles deben conservar la semántica de la señal; entre ellas se encuentran pequeñas variaciones de grosor, contraste, desplazamiento y escala dentro de límites conocidos. La Figura~\ref{fig:synthetic_examples} muestra una señal de cada familia.

\begin{figure}[H]
  \centering
  \safeincludegraphics[width=\textwidth]{results/fig_synthetic_examples.pdf}
  \caption{Ejemplos sintéticos por familia generadora utilizados para validar visualmente el simulador.}
  \label{fig:synthetic_examples}
\end{figure}

\section{Construcción del conjunto de datos}

Para construir el conjunto sintético se asigna una semilla propia a cada muestra y se genera la señal a partir de una de las familias descritas. Tras la aceptación, las muestras se distribuyen entre entrenamiento, validación y prueba. Las imágenes y su manifiesto se almacenan conjuntamente.

El manifiesto relaciona cada muestra con su imagen, partición, familia generadora, semilla, subtipo cuando corresponde, número de intento, índices internos y etiquetas binarias. Otros ficheros recogen el esquema de etiquetas y las frecuencias por partición, familia y combinación de clases. El entrenamiento puede leer directamente los manifiestos CSV, sin depender de la organización de las carpetas, mientras que los ficheros JSON facilitan la revisión del proceso y la elaboración de resultados.

\section{Comprobación de integración}

La configuración final consta de 100 pacientes sintéticos, cada uno con una imagen V1, sus pliegues LOOCV y los resúmenes de etiquetas. Veinte pacientes presentan los tres hallazgos; los 80 restantes son normales o contienen combinaciones parciales. Como cada configuración se repite con tres semillas, las matrices de confusión agregadas reúnen 300 predicciones.

Este conjunto constituye la base de la evaluación sintética. La prueba de integración abarca la generación de las señales, su aceptación mediante validadores, la partición por paciente y la escritura de los manifiestos. La coexistencia de varias combinaciones permite estudiar la evolución del rendimiento con \(K\) sin limitar la tarea a distinguir entre normalidad y un único patrón patológico.

La familia de origen y las etiquetas finales desempeñan funciones diferentes. La primera modifica los rangos de muestreo para favorecer una morfología, pero los validadores actúan después sobre la señal resultante. Así, una muestra generada desde la familia con elevación del ST puede activar también RBBB, y otra orientada a la inversión de T puede conservar una deflexión terminal asociada a otros hallazgos. La tarea queda definida por los rasgos detectados en la señal, no por el nombre de la familia inicial. La Figura~\ref{fig:label_distribution} muestra las combinaciones obtenidas.

\begin{figure}[H]
  \centering
  \safeincludegraphics[width=0.90\textwidth]{results/fig_label_distribution.pdf}
  \caption{Distribución de combinaciones de etiquetas derivadas en el conjunto sintético LOOCV.}
  \label{fig:label_distribution}
\end{figure}

\section{Riesgos de atajos visuales}

Un resultado alto no implica necesariamente que el modelo haya aprendido la morfología de interés. También podría aprovechar límites automáticos de los ejes, diferencias de antialiasing, márgenes ligados a la amplitud, cambios en la densidad de la cuadrícula, nombres de fichero o la propia estructura de carpetas.

Para reducir estos riesgos, las imágenes se leen desde el manifiesto sin proporcionar al modelo la familia de origen. Los límites visuales son constantes y los nombres de fichero solo contienen un contador. Como controles de sensibilidad se contemplan recortes de la zona del trazado y versiones sin el texto de los ejes. Una caída marcada al retirar elementos ajenos a la señal indicaría la necesidad de investigar posibles atajos antes de interpretar el rendimiento.

La atención espacial de la CNN se examina además mediante Grad-CAM \cite{selvaraju2017}. Este análisis no permite establecer relaciones causales; sirve para observar si las activaciones se concentran en QRS, ST y T o en elementos periféricos de la imagen.

\section{Limitaciones del simulador}

El simulador representa un latido aislado con una geometría semejante a la de una derivación precordial. No modela las doce derivaciones ni una transformación eléctrica común entre ellas. Tampoco incorpora variabilidad entre latidos, arritmias, intervalos dependientes de la frecuencia, ruido muscular, deriva de la línea base, interferencia de red, pérdida de electrodos o compresión de documentos.

La construcción mediante gaussianas produce curvas suaves y simplifica bordes que pueden ser más irregulares en un ECG real. A su vez, los validadores utilizan ventanas y umbrales diseñados exclusivamente para este dominio. Aplicarlos a una señal clínica como si fueran reglas diagnósticas carecería de justificación.

La familia Brugada simplifica una elevación del ST de morfología convexa asociada al patrón tipo 1. No representa la heterogeneidad del síndrome, la variación entre derivaciones ni su posible carácter intermitente. Se utiliza únicamente para estudiar la composición visual de los tres hallazgos.

Estas limitaciones delimitan el alcance del experimento. El dominio sintético permite comparar la eficiencia en el uso de ejemplos bajo condiciones controladas, pero no establece validez clínica. Esta última debe juzgarse a partir de la evaluación con la cohorte del HUCA y de una validación externa posterior.
